%%%%%%%%%%%%%%%%%%%%%%%%%%%%%%%%%
%%                             %%
%% Citovanie externých zdrojov %%
%%                             %%
%%%%%%%%%%%%%%%%%%%%%%%%%%%%%%%%%
\chapter{Citovanie externých zdrojov}\label{sec:citation}
V~súvislosti s~preberaním časti obsahu iných diel
rozoznávame dva pojmy. Sú to citát a~citácia.

Citát je doslovná reprodukcia prevzatého textu,
ktorý môžeme v~práci použiť dvomi spôsobmi.
Buď ako súčasť odseku textu,
alebo celý citovaný text vysádžeme v~samostatnom odseku.
V oboch prípadoch je zvykom citovaný text uzavrieť do úvodzoviek
a~zvýrazniť šikmým rezom písma.
Za citovaným textom uvedieme meno autora,
prípadne názov diela, rok a nasleduje číslo bibliografického
zdroja v~hranatých zátvorkách uvádzajúce poradie
v~zozname použitej literatúry v~závere práce.

Na citovanie v rámci odseku môžeme využiť makrá \verb|\uv| pre
správne slovenské úvodzovky a \verb|\emph|,
ktoré zabezpečí vytlačenie textu kurzívou.

Samostatne vysádzaný citát uzavrieme v \LaTeX-u do prostredia
\verb|\begin{quote}...| \verb|\end{quote}|.

Citácia je nepriamo prebraná a~prerozprávaná
časť citovanej práce.
Môže to byť vedecká myšlienka, odkaz na výsledky výskumu,
dôležitý poznatok, matematický vzťah a~podobne.
Schopnosť študenta pracovať s~literatúrou a~s~externými zdrojmi
je dôležitý moment pri hodnotení spôsobilosti uchádzača
o~vysokoškolský titul.
Tomuto aspektu práce treba preto venovať patričnú pozornosť.

Užitočný prehľad spôsobov citácií ponúkajú autorky
Beáta Bellérová a~Lucia Lichnerová v~článku v~časopise ITLib~\cite{Lichnerova2023Nove}.

\section{Odkazy na citované diela}
Vo vedecko-technických oblastiach, do ktorých patria aj študijné programy na našej fakulte, je zvykom používať numerický systém citovania. Jednotlivé zdroje sú očíslované podľa poradia výskytu v texte, pričom číslo zdroja uvádzame v hranatých zátvorkách.
Citovanie sa riadi technickou normou STN ISO 690: 2022 Informácie a dokumentácia: Návod na tvorbu bibliografických odkazov na informačné pramene a~ich citovanie~\cite{iso690}.

Systém \LaTeX\ umožňuje pracovať s citáciami niekoľkými spôsobmi.
Šablóna FEIstyle pracuje s balíčkom \verb|biblatex| na správu citácií. Jeho základom je externý databázový súbor \texttt{bibliography.bib},
ktorý sa nachádza v hlavnom priečinku projektu záverečnej práce.
Súbor obsahuje bibliografické záznamy citovaných diel
v špecifickom formáte.
Každý záznam začína jedinečným identifikátorom, ktorý zvyčajne
volíme tak, aby sme si ho jednoducho pamätali.
Ak totiž v texte chceme dielo citovať, napíšeme makro \verb|\cite{identifikator}| a~v~práci sa automaticky objaví poradové číslo citovaného diela v hranatej zátvorke.
Číslo citácii pridelí externý program \textsl{biber},
ktorý prejde celý dokument, identifikuje výskyty makra \verb|\cite|, zoradí citované diela podľa výskytu v~práci,
pridelí im čísla a vytvorí podklady na zoznam literatúry na konci práce.
Pri ďalšom spustení kompilátora \TeX-u dôjde k~nahradeniu makier
\verb|\cite| číslami v hranatých zátvorkách a makro šablóny
\verb|\FEIbibliography| vytvorí zoznam literatúry s nadpisom
prvej úrovne Literatúra.

Pri použití bibliografie odporúčame použiť dávkový súbor
\texttt{Makefile} alebo spustiť externý program \texttt{biber}
po prvej kompilácii a~potom skompilovať dokument ešte dvakrát.
Posledná kompilácia zabezpečí správne čísla strán v obsahu.
Online nástroj Overleaf robí všetky potrebné kroky v jednom behu,
kompiláciu tu stačí spustiť iba raz.

\section{Bibliografické záznamy}
V zozname použitej literatúry v~závere práce sa nachádzajú
podrobné záznamy použitých zdrojov.
Sú to najmä mená autorov, názov článku alebo číslo kapitoly
knihy, názov časopisu alebo knihy, vydavateľ, rok vydania,
strana, na ktorej sa nachádza citovaný článok a podobne.
V~prípade webových stránok je potrebné uviesť internetovú adresu
a~dátum, kedy sme informáciu zo stánky čerpali.
Dôležité je, aby boli informácie na základe týchto detailov
ľahko a jednoznačne dohľadateľné.
Norma STN ISO 690: 2022 tiež definuje formu takéhoto zoznamu~\cite{iso690}.

V~tejto publikácii formátujeme použitú literatúru
podľa štýlu iso-numeric \cite{Hoftich2022iso690} balíka \texttt{biblatex},
\footnote{\href{https://ctan.org/pkg/biblatex}{\texttt{ctan.org/pkg/biblatex}}}
ktorý do veľkej miery rešpektuje doteraz zaužívané zvyklosti
a~je navrhnutý tak,
aby spĺňal odporúčania normy aj v~slovenskom jazyku.

\subsection{Príklad záznamu v databázovom súbore .bib}\label{sec:citExample}
Súbor \texttt{bibliography.bib} je textový súbor, ktorý musí mať predpísanú štruktúru.
Môžeme ho vytvoriť ručne alebo použiť niektorý zo systémov na správu publikácií a~citácií ako sú JabRef, Mendeley, Zotero a podobne.

Nasledujúci bibliografický záznam
\begin{trivlist}
  \item STEINEROVÁ, J. Princípy formovania vzdelania
  v~informačnej vede. \textit{Pedagogická revue.} 2000, \textbf{2}(3), 8--16.
\end{trivlist}
vyzerá v zdrojovom súbore \texttt{bibliography.bib} takto:
\begin{verbatim}
@article{Steinerova2000Principy,
  author  = {J. Steinerová},
  journal = {Pedagogická revue},
  title   = {Princípy formovania vzdelania v informačnej vede},
  year    = {2000},
  number  = {3},
  pages   = {8--16},
  volume  = {2},
}
\end{verbatim}

Každý záznam začína symbolom \verb|@| a nasleduje typ publikácie definovaných v dokumentácii k balíčku Bib\LaTeX.
Okrem článku (\verb|article|) to môže byť aj kniha (\verb|book|),
príspevok v zborníku konferencie (\verb|inproceedings|),
webová stránka (\verb|online|),
správa (\verb|techreport|),
všeobecný záznam (\verb|misc|) a mnohé iné.

Prvý povinný údaj je jedinečný identifikátor, ktorý je ľubovoľný.
Odporúča sa však aby obsahoval priezvisko prvého autora, rok vydania a prvé slovo názvu.

Jednotlivé položky databázového záznamu sú oddelené čiarkou,
za posledným poľom už čiarka nie je.

Každý typ záznamu má iné požiadavky na polia.
Niektoré základné vlastnosti opíšeme v nasledujúcom texte.
Viac informácií možno získať v príslušnej dokumentácii~\cite{Hoftich2022iso690}.

\subsection{Prvky bibliografického záznamu}
Zoznam použitej literatúry obsahuje informácie o jednotlivých zdrojoch. Je potrebné mať na pamäti, aby záznamy boli jasné, stručné a jednoznačne identifikovateľné. Bežne zorientovaný čitateľ práce, ktorý sa pohybuje v okruhu tém práce by mal byť schopný identifikovať jednotlivé diela na prvé pozretie. Prípadný záujemca o detailné štúdium problematiky by si mal vedieť na základe záznamu vyhľadať citovanú publikáciu buď na internete alebo v knižnici. Týmito pravidlami sa riadime pri vytváraní zoznamu použitej literatúry.

Uvedieme niekoľko jednoduchých a~najčastejších príkladov bežného
spôsobu citovania.
Snaha je, aby sme sa prirodzene naučili formy uvádzania
bibliografických záznamov v~našej oblasti vedy a techniky
a~dokázali vytvoriť správnu databázu zdrojov v systéme Bib\LaTeX.

\subsection*{Mená tvorcov}
V zozname literatúry majú formu: PRIEZVISKO, Meno alebo PRIEZVISKO, M. a~navzájom sú oddelené bodkočiarkou.
Za zoznamom autorov nasleduje bodka.

\subsection*{\normalsize Pole v .bib súbore}
\begin{verbatim}
author = {Meno1 Priezvisko1 and Meno2 Priezvisko2 and Meno3 Prizvisko3}
\end{verbatim}
Autorov zapisujeme ako \verb|Meno Priezvisko| a~oddeľujeme ich kľúčovým slovom and.
Môžeme tiež použiť obrátený tvar \verb|Priezvisko, Meno and|\dots s~čiarkou medzi priezviskom a~menom.

\subsection*{Názov}
Názov článku píšeme normálnym písmom.
Názov nosného informačného zdroja (kniha, časopis, zborník) píšeme kurzívou.

\subsection*{\normalsize Polia v .bib súbore}
\begin{verbatim}
title = {}
booktitle = {}
\end{verbatim}

\subsection*{Dátum}
Pri každej publikácii musí byť uvedený aspoň rok vydania
vo formáte YYYY, t. j. 2024.
V prípade online zdrojov treba uvádzať aj presný dátum,
kedy sme zdroj použili.
Vkladáme ho do hranatých zátvoriek s~kľúčovým slovom cit.
Odporúčame použiť formát [cit. YYYY-MM-DD].
Ak sme teda informáciu z webovej stránky získali dňa
14. mája 2023,
napíšeme to takto: [cit. 2023-05-14].
Ide o medzinárodne akceptovaný a~zrozumiteľný tvar zápisu dátumu.

\subsection*{\normalsize Polia v .bib súbore}
\begin{verbatim}
year = {2024}
date = {2024-12-31}
\end{verbatim}

\subsection*{Dostupnosť}
Myslí sa tým dostupnosť na internete najmä prostredníctvom webu.
Uvádzame buď úplnú webovú adresu alebo DOI číslo.
Nepíšeme oba údaje, preferujeme DOI.
Tento údaj sa nachádza zväčša na konci záznamu za kľúčovým výrazom Dostupné na.

\subsection*{\normalsize Polia v .bib súbore}
\begin{verbatim}
url = {}
doi = {}
eprint = {}
\end{verbatim}
Ak uvedieme viacero polí, Bib\LaTeX\ zväčša vyberie iba jedno
podľa preferencií špecifikovaných v~definičnom súbore.

\subsection*{Ročník alebo zväzok a číslo časopisu}
Vedecké periodiká vychádzajú v tzv. zväzkoch (angl. \foreignlanguage{english}{\it volume}).
Zväzky sa môžu deliť na čísla
(angl. \foreignlanguage{english}{\it issue} alebo
\foreignlanguage{english}{\it number}).
Zväzok a číslo zapisujeme buď pomocou skratiek vol. a~iss.
(prípadne no. podľa skutočného členenia jednotlivých vydaní konkrétneho časopisu),
alebo preferujeme zaužívanú skrátenú formu,
pri ktorej nepoužijeme slovné skratky,
ale píšeme len čísla,
pričom zväzok vytlačíme tučným rezom (bold)
tesne nasledovaný číslom časopisu normálnym písmom
v~okrúhlych zátvorkách.
Napríklad zväzok 15, číslo 8 môžeme zapísať
ako vol. 15, iss. 8 alebo \textbf{15}(8).

\subsection*{\normalsize Polia v .bib súbore}
\begin{verbatim}
volume = {}
number = {}
\end{verbatim}

\section{Článok v odbornom periodiku}
Záznam musí obsahovať mená autorov (tvorcov), názov článku, názov časopisu, dátum alebo iba rok vydania, ročník alebo zväzok (volume) a číslo vo zväzku (issue), číslo prvej strany, prípadne rozsah strán, na ktorých sa nachádza.

\subsection*{\normalsize Nepovinné údaje}
Ak je článok dostupný online, je dobré, ak sa v zápise objaví aj webová adresa digitálnej verzie článku. Vedecké publikácie majú pridelený jedinečný kód, tzv. DOI (Digital Object Identifier – identifikátor digitálneho objektu), ktorý tiež môže byť súčasťou bibliografického zápisu. Ďalšou nepovinnou časťou je International Standard Serial Number (ISSN), teda medzinárodné štandardné sériové číslo, ktoré sa prideľuje časopisom.

\subsection*{\normalsize Printový časopis}
\begin{trivlist}
\item TVORCOVIA. Názov článku. \textit{Názov periodika.} Rok vydania, \textbf{zväzok}(časť), rozsah strán. Dostupnosť.
\end{trivlist}

\subsection*{\normalsize Online časopis}
\begin{trivlist}
\item TVORCOVIA. Názov článku. \textit{Názov periodika} [online]. Rok vydania, \textbf{zväzok}(časť), rozsah strán (ak je k dispozícii). [cit. yyyy-mm-dd]. Dostupnosť.
\end{trivlist}

\subsection*{\normalsize Príklady}
Prvý príklad je z~časti
\ref{sec:citExample} odkazuje na článok s názvom Princípy formovania vzdelania v informačnej vede od autorky Jely Steinerovej, ktorý vyšiel v roku 2000 v 3. čísle 2. ročníka časopisu Pedagogická revue. Článok sa v časopise nachádza na stranách 8 až 16.

V druhom príklade skratka et. al. (skratka latinského slovného spojenia \textit{et alii}, vo význame a~iní) za menom autora znamená, že článok má viacero autorov (bežne viac než 5) a nemusíme ich všetkých uvádzať.
V slovenských textoch sa môže použiť skratka a~kol. (a~kolektív).

Tretí príklad je príklad online časopisu, ktorý možno nájsť na internete. Príznak [online] a~dátum citovania [cit. 2023-01-10] sa objaví v zozname literatúry, keď v súbore \verb|bibliography.bib| použijeme pole \verb|urldate|.

\begin{enumerate}
\item STEINEROVÁ, J. Princípy formovania vzdelania v informačnej vede. \textit{Pedagogická revue}. 2000, \textbf{2}(3), 8–16.
\item BEŇAČKA, J. et al. A better cosine approximate solution to pendulum equation. \textit{International Journal of Mathematical Education in Science and Technology}, 2009, \textbf{40}(2), 206–215.
\item HOGGAN, D. B. Challenges, Strategies, and Tools for Research Scientists. Electronic Journal of \textit{Academic and Special Librarianship} [online]. 2009, \textbf{3}(3). [cit. 2023-01-10]. Dostupné z: \texttt{http://southernlibrarianship.icaap.org/content/ v03n03/Hoggan\_d01.htm}. ISSN 1525-321X.
\end{enumerate}

\subsection*{\normalsize Zápis v zdrojovom .bib súbore}
\begin{verbatim}
@article{Steinerova2000Principy,
  author  = {J. Steinerová},
  journal = {Pedagogická revue},
  title   = {Princípy formovania vzdelania v informačnej vede},
  year    = {2000},
  number  = {3},
  pages   = {8--16},
  volume  = {2},
}

@article{Benacka2009Abetter,
  author  = {Ján Benačka and others},
  title   = {A better cosine approximate solution to pendulum
             equation},
  journal = {International Journal of Mathematical Education in Science
             and Technology},
  volume  = {40},
  number  = {2},
  pages   = {307--308},
  year    = {2009},
  doi     = {10.1080/00207390802419594},
}

@article{Hoggan2002Challenges,
  author  = {Danielle Bodrero Hoggan},
  journal = {Electronic Journal of Academic and Special Librarianship},
  title   = {Challenges, Strategies, and Tools for Research Scientists:
             Using Web-Based Information Resources},
  year    = {2002},
  number  = {3},
  volume  = {3},
  url     = {https://southernlibrarianship.icaap.org/content/v03n03/
             Hoggan_d01.htm},
  urldate = {2023-01-10},
}
\end{verbatim}

\section{Monografia a~kniha}
\begin{trivlist}
\item TVORCOVIA. Názov knihy. Vydanie. Mesto: Vydavateľ, Rok vydania. ISBN.
\end{trivlist}

\subsection*{\normalsize Príklady}
\begin{enumerate}
\item OBERT, V. \textit{Návraty a odkazy}. Nitra: Univerzita Konštantína Filozofa, 2006. ISBN 80-8094-046-0.

\item TIMKO, J.; SIEKEL, P.; TURŇA, J. Geneticky modifikované organizmy. Bratislava: Veda, 2004. ISBN 80-224-0834-4.
\end{enumerate}

\subsection*{\normalsize Zápis v zdrojovom .bib súbore}
\begin{verbatim}
@book{Obert2006Navraty,
  author    = {Viliam Obert},
  publisher = {Univerzita Konštantína Filozofa},
  title     = {Návraty a odkazy},
  year      = {2006},
  address   = {Nitra},
  isbn      = {80-8094-046-0},
}

@book{Timko2004Geneticky,
  author    = {Jozef Timko and Peter Siekel and Ján Turňa},
  publisher = {Veda},
  title     = {Geneticky modifikované organizmy},
  year      = {2004},
  isbn      = {80-224-0834-4},
}
\end{verbatim}

\section{Záverečná a~vedecko-kvalifikačná práca, správa}
AUTOR. \textit{Názov práce}. Mesto, Rok vypracovania. Typ práce. Inštitúcia

\subsection*{\normalsize Príklady}
\begin{enumerate}
  \item MIKULÁŠIKOVÁ, M. \textit{Didaktické pomôcky pre praktickú výučbu na hodinách výtvarnej výchovy pre 2. stupeň základných škôl.} Nitra, 1999. Diplomová práca. Univerzita Konštantína Filozofa

  \item BAUMGARTNER, J. et al. \textit{Ochrana a udržiavanie genofondu zvierat, šľachtenie zvierat.} Nitra, 1998. Výskumná správa. VÚŽV.
\end{enumerate}

\subsection*{\normalsize Zápis v zdrojovom .bib súbore}
\begin{verbatim}
@thesis{Mikulasikova1999Didakticke,
  author  = {M. Mikulášiková},
  title   = {Didaktické pomôcky pre praktickú výučbu na hodinách
             výtvarnej výchovy pre 2. stupeň základných škôl},
  address = {Nitra},
  year    = {1999},
  type    = {Diplomová práca},
  school  = {Univerzita Konštantína Filozofa},
}

@report{Baumgarntner1998Ochrana,
  author      = {J. Baumgartner and others},
  institution = {VÚŽV},
  title       = {Ochrana a udržiavanie genofondu zvierat, šľachtenie
                 zvierat},
  address     = {Nitra},
  year        = {1998},
  type        = {Výskumná správa},
}
\end{verbatim}

\section{Príspevok v zborníku konferencie, časť knihy}
Ak citujeme iba jednu kapitolu knihy,
prípadne ide o~príspevok z~neperiodického zborníka,
uvedieme za názvom príspevku (kapitoly)
kľúčové slovo In:, po dvojbodke nasleduje zoznam editorov zborníka a názov zborníka alebo knihy:
\begin{trivlist}
  \item TVORCOVIA. Názov príspevku. In: EDITORI (ed.). \textit{Názov zborníka.} Mesto: Vydavateľ, Rok vydania, zväzok, s. rozsah strán, zv. číslo. ISBN. ISSN. Dostupnosť.
\end{trivlist}

\subsection*{\normalsize Príklady}
\begin{enumerate}
  \item ZEMÁNEK, P. The machines for ``green works'' in vineyards and their economical evaluation. In: \textit{9th International Conference: proceedings. Vol. 2. Fruit Growing and viticulture.} Lednice: Mendel University of Agriculture a Forestry, 2001, zv. 2, s. 262–268. ISBN 80-7157-524-0.

  \item CHLPÍK, J.; KURTULÍK, M.; KOTOROVÁ, S.; CIRÁK, J. Spectroscopic ellipsometry of Au nanoparticles layers. In: SITEK, J.; VAJDA, J.; JAMNICKÝ, I. (ed.). \textit{AIP Conference Proceedings.} 2024, zv. 3054, s. 080004. Č. 1. ISSN 0094-243X. Dostupné z doi: 10.1063/5.0187526.
\end{enumerate}

\subsection*{\normalsize Zápis v zdrojovom .bib súbore}
\begin{verbatim}
@InProceedings{Zemanek2001TheMachines,
  author    = {P. Zemánek},
  booktitle = {9th International Conference: proceedings. Vol. 2. Fruit
               Growing and viticulture},
  title     = {The machines for ``green works'' in vineyards and their
               economical evaluation},
  volume    = {2},
  year      = {2001},
  pages     = {262--268},
  publisher = {Mendel University of Agriculture and Forestry},
  address   = {Lednice},
  isbn      = {80-7157- 524-0},
}

@InProceedings{ChlpikSpectroscopic2024,
  author    = {Chlpík, Juraj and Kurtulík, Matej and Kotorová, Soňa
               and Cirák, Július},
  date      = {2024-01},
  title     = {Spectroscopic ellipsometry of Au nanoparticles
               layers},
  editor    = {Jozef Sitek and Ján Vajda and Igor Jamnický},
  doi       = {10.1063/5.0187526},
  number    = {1},
  pages     = {080004},
  volume    = {3054},
  issn      = {0094-243X},
  booktitle = {AIP Conference Proceedings},
}
\end{verbatim}

\section{Webová stránka, sociálna sieť, video}
Pri multimediálnom online obsahu často nepoznáme autora,
prípadne je komplikované zistiť názov webovej stránky.
Snažíme sa teda zahrnúť čo najviac jednoznačných informácií,
najmä webovú adresu, dátum publikovania a dátum citovania.
\begin{trivlist}
  \item TVORCOVIA. \textit{Názov obsahu} [online]. Platforma, dátum publikovania [cit. dátum citovania]. Dostupnosť.
\end{trivlist}

\subsection*{\normalsize Príklady}
\begin{enumerate}
  \item VALKO, P. \textit{M31, M32 a jeden Starlink k tomu} [online]. Facebook, 2024-08-30 [cit. 2024-09-06]. Dostupné z: \texttt{https://www.facebook.com/share/p/S93xZoz9Rbnh
  dQ7M}.

  \item \textit{Why lenses can't make perfect images} [online]. Youtube, 2017-10-18 [cit. 2024-08-25]. Dostupné z: \texttt{https://youtu.be/DDoryfCXxPI?si=hC5Kuuf4p3OxOOsm}.
\end{enumerate}

\subsection*{\normalsize Zápis v zdrojovom .bib súbore}
\begin{verbatim}
@online{Valko2024M31,
  author       = {Pavol Valko},
  title        = {31, M32 a jeden Starlink k tomu},
  howpublished = {online},
  date         = {2025-08-30T00:03:00},
  url          = {https://www.facebook.com/share/p/S93xZoz9RbnhdQ7M},
  urldate      = {2024-09-06},
  publisher    = {Facebook},
}

@online{WhyLenses2017,
  title        = {Why lenses can’t make perfect images},
  howpublished = {online},
  date         = {2025-10-18},
  url          = {https://youtu.be/DDoryfCXxPI?si=hC5Kuuf4p3OxOOsm},
  urldate      = {2024-08-25},
  publisher    = {Youtube},
}
\end{verbatim}

\section{Ako citovať technické normy}
K normám nemusíme uvádzať autorov, ak ide o slovenskú normu STN, netreba ani vydavateľa. Dôležité je číslo normy a rok vydania.

\begin{trivlist}
\item \textit{Číslo normy. Názov: Podnázov.} Mesto: Vydavateľ. Rok vydania.
\end{trivlist}

\subsection*{\normalsize Príklad}
\begin{itemize}
  \item \textit{STN ISO 690: 2023. Informácie a~dokumentácia: Návod na tvorbu bibliografických odkazov na informačné pramene a ich citovanie.} Bratislava: Slovenský ústav technickej normalizácie, 2023.
\end{itemize}

\subsection*{\normalsize Zápis v zdrojovom .bib súbore}
\begin{verbatim}
@report{iso690,
  title     = {STN ISO 690: 2022. Informácie a~dokumentácia},
  subtitle  = {Návod na tvorbu bibliografických odkazov na informačné
               pramene a~ich citovanie},
  address   = {Bratislava},
  publisher = {Slovenský ústav technickej normalizácie},
  year      = {2022},
}
\end{verbatim}

% LocalWords:  zostručnenie Bib vý nok výk le ok ax bx ds dt dx Eq CoulombVec
% LocalWords:  eq px lrrrr FWHM Full Width rrrr kA Technology HOGGAN Challenges
% LocalWords:  Strategies Tools Research Scientists Electronic Academic Special
% LocalWords:  Librarianship Hoggan OBERT TIMKO SIEKEL MIKULÁŠIKOVÁ BAUMGARTNER
% LocalWords:  VÚŽV ed ZEMÁNEK machines green works vineyards their economical
% LocalWords:  evaluation th Conference proceedings Fruit Growing viticulture
% LocalWords:  Mendel University Agriculture Forestry CHLPÍK KURTULÍK KOTOROVÁ
% LocalWords:  CIRÁK Spectroscopic ellipsometry nanoparticles layers SITEK doi
% LocalWords:  JAMNICKÝ VALKO Starlink Facebook dQ Why lenses can perfect
% LocalWords:  images Youtube
